% Offline Crop Advisor: Technical Report
% Bharat AI-SoC Challenge 2026-27, PS1
% Build:  pdflatex technical_report.tex   (run twice for the table of contents)

\documentclass[11pt,a4paper]{article}

\usepackage[T1]{fontenc}
\usepackage[utf8]{inputenc}
\usepackage{lmodern}
\usepackage[margin=2.4cm]{geometry}
\usepackage{microtype}
\usepackage{booktabs}
\usepackage{tabularx}
\usepackage{array}
\usepackage{enumitem}
\usepackage{xcolor}
\usepackage{titlesec}
\usepackage{fancyhdr}
\usepackage{fancyvrb}
\usepackage{float}         % [H]: keep tables exactly where they are written
\usepackage{caption}
\usepackage{xurl}          % let long URLs break anywhere
\definecolor{linkblue}{RGB}{13,71,161}
\usepackage[colorlinks=true,urlcolor=linkblue,linkcolor=linkblue,
            citecolor=linkblue,breaklinks=true]{hyperref}
\usepackage{bookmark}      % stable PDF outlines in one extra run

% ---------------------------------------------------------------- styling ---
\definecolor{accent}{RGB}{27,94,32}

\titleformat{\section}{\large\bfseries\color{accent}}{\thesection.}{0.6em}{}
\titlespacing*{\section}{0pt}{1.4\baselineskip}{0.5\baselineskip}
\titleformat{\subsection}{\normalsize\bfseries}{\thesubsection}{0.6em}{}
\titlespacing*{\subsection}{0pt}{1.0\baselineskip}{0.35\baselineskip}

\setlist[itemize]{leftmargin=1.3em,topsep=0.3em,itemsep=0.25em}
\setlist[enumerate]{leftmargin=1.6em,topsep=0.3em,itemsep=0.25em}

\pagestyle{fancy}
\fancyhf{}
\renewcommand{\headrulewidth}{0.4pt}
\lhead{\small Offline Crop Advisor --- Technical Report}
\rhead{\small\thepage}

\renewcommand{\arraystretch}{1.25}
\newcolumntype{R}{>{\raggedleft\arraybackslash}X}
\newcolumntype{L}{>{\raggedright\arraybackslash}X}
\newcommand{\hdr}[1]{\textbf{#1}}

% Narrow helper for numeric columns
\newcolumntype{N}{>{\raggedleft\arraybackslash}p{1.45cm}}

\begin{document}

% ------------------------------------------------------------------ title ---
\begin{center}
  {\LARGE\bfseries Offline Crop Advisor: Technical Report}\\[0.9em]
  {\large Bharat AI-SoC Challenge 2026--27, PS1}\\[0.35em]
  {\normalsize Offline crop disease and pest detection and advisory for farmers}\\[0.9em]
  {\normalsize Phase 1 (Open Qualifying) submission, 23 November 2026}\\[0.5em]
  {\normalsize\bfseries Team: Aishani Roy, Vijoya Guha Thakurta}\\[0.6em]
  {\small Source code: \url{https://github.com/aishani0102/offline-crop-advisor}}\\[0.3em]
  {\small Demo video: \url{https://drive.google.com/file/d/1YLU4_h4lRAjxv5cIshW1-GhEhf6i896o/view?usp=drive_link}}
\end{center}

\vspace{0.6em}
\hrule
\vspace{1.2em}

% ---------------------------------------------------------------- summary ---
\section{Summary}

Offline Crop Advisor is an Android app that tells a farmer what is wrong with a
tomato or potato leaf and what to do about it, with no internet connection at any
point. The farmer picks the crop, takes a photo, answers three short questions,
and gets:

\begin{itemize}
  \item the disease or pest (10 classes, including one pest), with a confidence score,
  \item how much of the leaf is affected (mild, moderate or severe),
  \item advice in their language (English, Hindi, Bengali, Marathi, Tamil or
        Telugu), written by \textbf{Llama 3.2 1B running on the phone with
        ExecuTorch}, grounded in a knowledge base of Indian extension guidance,
        with doses, a pre-harvest interval and a KVK / Kisan Call Centre referral.
\end{itemize}

Everything runs on an Arm phone (Redmi 12 5G, Snapdragon 4 Gen 2). Tested in
airplane mode: 10 of 10 full runs, from photo to Llama advice, finished offline
and passed every safety check.

\begin{table}[H]
\centering
\begin{tabularx}{\textwidth}{@{}L L@{}}
\toprule
\hdr{Headline} & \hdr{Result} \\
\midrule
Leaf classifier (MobileNetV3-Large, INT8, 4.5\,MB) & Field test: 93.2\% accuracy, macro F1 0.893 (1,020 field photos). Lab test: 98.5\% \\
Classifier time on the phone & 70 to 300\,ms per photo \\
Llama 3.2 1B (SpinQuant INT4, 1.1\,GB) on the phone & about 12 tokens/s; first token 1.3\,s (short prompt) to 8\,s (full grounded prompt) \\
Advice shown to the farmer that passed every safety check & 60 of 60 test cases \\
Peak app memory with Llama loaded & about 1.55\,GB \\
Energy (Android estimate, CPU only) & about 3.2\,mAh per diagnosis, about 1,360 diagnoses per charge \\
\bottomrule
\end{tabularx}
\end{table}

% -------------------------------------------------------- problem & scope ---
\section{Problem and scope}

Small farmers lose a large share of tomato and potato yield to a handful of
diseases, mostly blights. Expert help (KVKs, extension officers) is far away and
mobile data is patchy in fields. The challenge asks for detection, severity,
treatment advice, regional relevance and local languages, all on-device.

Scope decisions made up front:

\begin{itemize}
  \item \textbf{Two crops, ten classes.} Tomato: healthy, early blight, late
        blight, leaf mold, Septoria leaf spot, yellow leaf curl virus, leaf miner
        (pest). Potato: healthy, early blight, late blight. These cover the most
        common problems in Indian fields while keeping enough field photos per class.
  \item \textbf{Android phone as the target.} Farmers already own Arm Android
        phones. The same models run on a Raspberry Pi 5, but the phone is the
        realistic device.
  \item \textbf{The farmer chooses the crop.} This removes tomato and potato
        confusion and halves the classes the model has to pick from.
  \item \textbf{Advice must never invent chemicals or doses.} The language model
        rewrites facts from a knowledge base; a rule-based checker rejects
        anything not in it.
\end{itemize}

% -------------------------------------------------------- system overview ---
\section{System overview}

\begin{Verbatim}[fontsize=\footnotesize,frame=single,framesep=6pt,rulecolor=\color{gray}]
Photo
  -> Quality gate        too dark / too blurry / no leaf  -> ask for a retake
  -> Leaf classifier     MobileNetV3-Large INT8 (ExecuTorch), chosen crop only
                         not this crop / confidence < 0.56 -> "not sure", refer to KVK
  -> Severity            colour segmentation -> mild / moderate / severe
  -> 3 questions         rain, crop stage, already sprayed (+ month, state)
  -> KB facts            doses, pre-harvest interval, referral
  -> Llama 3.2 1B        rewrites the facts in the farmer's language (on device)
  -> Safety check        pass -> Llama text    fail -> KB template text
\end{Verbatim}

\begin{enumerate}
  \item \textbf{Quality gate.} Rejects photos that are too dark, too blurry, or
        have almost no leaf-coloured pixels, and asks for a retake. No diagnosis
        is shown for a rejected photo.
  \item \textbf{Leaf classifier.} MobileNetV3-Large, INT8, run with ExecuTorch
        (XNNPACK). Probabilities are temperature-calibrated and restricted to the
        chosen crop. If the photo does not look like the chosen crop, or
        confidence is below 0.56, the app says ``not sure'' and refers the farmer
        to the KVK instead of naming a treatment (section~5).
  \item \textbf{Severity.} Share of the leaf that is diseased, from a colour
        segmentation of leaf and lesion pixels, banded as mild ($<10\%$),
        moderate (10 to 30\%) or severe ($>30\%$). Only run for disease and pest labels.
  \item \textbf{Three questions.} Recent rain, crop stage and whether the crop was
        already sprayed. With the month and state, they order the advice (for
        example, no spraying close to harvest).
  \item \textbf{Advice.} The knowledge base gives the facts for the case. Llama
        3.2 1B rewrites them in plain language. A safety checker verifies the
        output; if it fails, the farmer sees the knowledge base template text,
        which is always safe.
\end{enumerate}

Only one model is in memory at a time apart from the small classifier, which
keeps the app within a 2\,GB free-RAM budget.

% ------------------------------------------------------------------- data ---
\section{Data}

Lab photos teach a model what a disease looks like; field photos teach it what a
real phone photo looks like. We combined one large lab dataset with six field
datasets so that every class has field photos, and we report accuracy only on
field photos the model never saw. More detail, figures and sample images are in
the separate Dataset Study (\texttt{Dataset\_Study.pdf}). Every number here comes
from the build manifest and is reproduced by
\texttt{detection/data/build\_dataset.py} with a fixed seed (split fingerprint
\texttt{09129439436a1ece}).

\subsection{Sources}

\begin{table}[H]
\centering
\small
\begin{tabularx}{\textwidth}{@{}l L L l@{}}
\toprule
\hdr{Source} & \hdr{Photos} & \hdr{Origin} & \hdr{Licence} \\
\midrule
PlantVillage & Lab, single leaf on a plain background & Penn State; Hughes and Salathe 2015 [1] & CC BY-SA 3.0 \\
PlantDoc & Field and web photos, cluttered backgrounds & IIT Gandhinagar; Singh et al.\ 2020 [2] & CC BY 4.0 \\
PlantSeg & Field and web photos, with lesion masks & Univ.\ of Queensland; Wei et al.\ 2024 [3] & CC BY-NC 4.0 \\
Tomato-Village & Field photos, Jodhpur and Jaipur, Rajasthan, India & Gehlot et al.\ 2023 [4] & No licence file \\
Potato, Indonesia & Field photos, Central Java & Shabrina et al.\ 2024 [5] & CC BY 4.0 \\
Potato, Ethiopia & Field photos, Holeta & Mendeley Data [6] & CC BY 4.0 \\
Potato, Punjab (PLD) & Field photos, Central Punjab, Pakistan, labelled with plant pathologists & Rashid et al.\ 2021 [7] & ODbL / DbCL \\
\bottomrule
\end{tabularx}
\end{table}

Tomato-Village has no licence file; it is downloaded by script and never
re-hosted. Left out on purpose: the Indonesia ``Fungi'' folder (several fungal
diseases under one label), classes outside our list, the other PlantVillage
tomato classes, and Kaggle re-uploads that are augmented copies of PlantVillage
(they would leak near-identical images into the test set).

\subsection{Cleaning}

Every image gets a perceptual hash; two images within 4 bits (out of 64) count as
the same photo, across all sources. When copies are found, the test copy is kept
and the training copy removed, so no test photo has a near-copy in training.

\begin{table}[H]
\centering
\begin{tabular}{@{}lr@{}}
\toprule
\hdr{Step} & \hdr{Images} \\
\midrule
Found in our 10 classes & 23,903 \\
Unreadable files & 0 \\
Near-duplicates removed & 795 \\
Removed because copies had different labels & 175 \\
\midrule
Left after cleaning & \textbf{22,933} \\
\bottomrule
\end{tabular}
\end{table}

Of the 795 near-duplicates, 524 were inside Tomato-Village.

\subsection{Balancing}

PlantVillage alone is badly skewed (5,357 lab photos of yellow leaf curl virus
against 152 of healthy potato). We balanced in three steps:

\begin{itemize}
  \item \textbf{Field photos first:} every field photo of a class goes into
        training, up to 700.
  \item \textbf{Lab photos fill the gap} up to 1,000 per class, which stops the
        large lab classes from swamping the rest.
  \item \textbf{The sampler evens out the rest:} each class is drawn about equally
        often per epoch, and a field photo is twice as likely to be drawn as a lab photo.
\end{itemize}

Training augmentation is built for field conditions: random crops, rotation,
colour and brightness changes, blur and blanked-out patches.

\subsection{Splits}

\begin{itemize}
  \item \textbf{Test, field only (1,020):} the official test folders of PlantDoc,
        PlantSeg, Tomato-Village and PLD, plus a 20\% hold-out of the Indonesia
        and Ethiopia sets. Never used in training. This is the headline accuracy.
  \item \textbf{Test, lab only (1,135):} PlantVillage photos held out from
        training, to show the gap between lab and field.
  \item \textbf{Validation (1,800):} used to pick the best epoch, calibrate
        confidence and choose the confidence threshold. Lab photos are capped at
        100 per class so field photos dominate.
  \item \textbf{Train (9,421):} everything else, after balancing; 46\% field photos.
\end{itemize}

\begin{table}[H]
\centering
\small
\caption*{\textbf{Images per source and split}}
\begin{tabular}{@{}lrrrrrr@{}}
\toprule
\hdr{Source} & \hdr{Found} & \hdr{Train} & \hdr{Val.} & \hdr{Test (field)} & \hdr{Test (lab)} & \hdr{Kept} \\
\midrule
PlantVillage & 14,732 & 5,115 & 810 & 0 & 1,135 & 7,060 \\
PlantDoc & 801 & 588 & 65 & 61 & 0 & 714 \\
PlantSeg & 680 & 392 & 64 & 122 & 0 & 578 \\
Tomato-Village & 2,640 & 1,458 & 382 & 257 & 0 & 2,097 \\
Potato, Indonesia & 548 & 354 & 50 & 101 & 0 & 505 \\
Potato, Ethiopia & 430 & 288 & 41 & 82 & 0 & 411 \\
Potato, Punjab (PLD) & 4,072 & 1,226 & 388 & 397 & 0 & 2,011 \\
\midrule
\textbf{Total} & \textbf{23,903} & \textbf{9,421} & \textbf{1,800} & \textbf{1,020} & \textbf{1,135} & \textbf{13,376} \\
\bottomrule
\end{tabular}
\end{table}

\begin{table}[H]
\centering
\small
\caption*{\textbf{Images per class and split}}
\begin{tabular}{@{}lrrrrr@{}}
\toprule
\hdr{Class} & \hdr{Train} & \hdr{of which field} & \hdr{Val.} & \hdr{Test (field)} & \hdr{Test (lab)} \\
\midrule
Tomato, healthy & 1,000 & 198 (20\%) & 148 & 30 & 158 \\
Tomato, early blight & 1,000 & 348 (35\%) & 172 & 85 & 100 \\
Tomato, late blight & 1,000 & 555 (56\%) & 207 & 113 & 190 \\
Tomato, leaf mold & 899 & 138 (15\%) & 115 & 24 & 95 \\
Tomato, Septoria leaf spot & 1,000 & 169 (17\%) & 125 & 27 & 177 \\
Tomato, yellow leaf curl virus & 1,000 & 98 (10\%) & 109 & 21 & 200 \\
Tomato, leaf miner (pest) & 700 & 700 (100\%) & 202 & 104 & 0 \\
Potato, healthy & 822 & 700 (85\%) & 169 & 209 & 15 \\
Potato, early blight & 1,000 & 700 (70\%) & 261 & 172 & 100 \\
Potato, late blight & 1,000 & 700 (70\%) & 292 & 235 & 100 \\
\midrule
\textbf{Total} & \textbf{9,421} & \textbf{4,306 (46\%)} & \textbf{1,800} & \textbf{1,020} & \textbf{1,135} \\
\bottomrule
\end{tabular}
\end{table}

The four tomato classes with 21 to 30 field test images (leaf mold, Septoria,
leaf curl virus, healthy) give noisy per-class scores and are reported with that
caveat.

% -------------------------------------------------------- detection model ---
\section{Detection model}

\begin{table}[H]
\centering
\begin{tabularx}{\textwidth}{@{}l L@{}}
\toprule
\hdr{Item} & \hdr{Choice} \\
\midrule
Architecture & MobileNetV3-Large, ImageNet pre-trained \\
Input & Shorter side resized to 256, centre crop 224, ImageNet normalisation \\
Training & 3 epochs head only, then 22 epochs full network; AdamW; class-balanced sampler with field photos weighted $2\times$; field-style augmentation (crops, rotation, colour, blur, occlusion) \\
Calibration & Temperature scaling on validation ($T = 0.687$) \\
Export & ExecuTorch \texttt{.pte}, XNNPACK; INT8 post-training quantisation calibrated on field and lab photos \\
\bottomrule
\end{tabularx}
\end{table}

\subsection{Results (INT8, the model in the app)}

\begin{table}[H]
\centering
\begin{tabular}{@{}lrrr@{}}
\toprule
\hdr{Test set} & \hdr{Images} & \hdr{Accuracy} & \hdr{Macro F1} \\
\midrule
Field photos (6 sources) & 1,020 & 0.932 & 0.893 \\
Lab photos (PlantVillage) & 1,135 & 0.985 & 0.986 \\
\bottomrule
\end{tabular}
\end{table}

FP32 before quantisation: field accuracy 0.953, macro F1 0.919. INT8 costs 2.6
points of field macro F1 and makes the model $4\times$ smaller (4.5\,MB).

\textbf{Per source (field test, INT8).}

\begin{table}[H]
\centering
\begin{tabular}{@{}lrrr@{}}
\toprule
\hdr{Source} & \hdr{Images} & \hdr{Accuracy} & \hdr{Macro F1} \\
\midrule
PlantDoc (web) & 61 & 0.705 & 0.679 \\
PlantSeg (field and web) & 122 & 0.795 & 0.799 \\
Tomato-Village (Rajasthan) & 257 & 0.938 & 0.929 \\
Potato, Indonesia & 101 & 0.980 & 0.979 \\
Potato, Ethiopia & 82 & 0.988 & 0.978 \\
Potato, Punjab (PLD) & 397 & 0.982 & 0.984 \\
\bottomrule
\end{tabular}
\end{table}

Single-site sets, where every photo was taken by one team in one place, score 94
to 99\%. The web-scraped sets, which mix many cameras and places, score 70 to
80\%; that is the better guess for photos from a new farm.

\textbf{Phone vs PC.} The app's preprocessing was checked against the Python
pipeline: on 20 healthy potato photos and 12 PLD photos the phone gave the same
labels as the PC, with confidences within 0.01.

\subsection{Why MobileNetV3-Large and not MobileNetV3-Small}

We also trained MobileNetV3-Small and exported it the same way (same quantiser,
same 256 calibration images).

\begin{table}[H]
\centering
\begin{tabular}{@{}lrrr@{}}
\toprule
\hdr{Model} & \hdr{Parameters} & \hdr{INT8 file} & \hdr{INT8 latency (laptop CPU, median)} \\
\midrule
MobileNetV3-Large & 4.2\,M & 4.5\,MB & 7.9\,ms \\
MobileNetV3-Small & 1.5\,M & 1.7\,MB & 5.5\,ms \\
\bottomrule
\end{tabular}
\end{table}

Small is $2.6\times$ smaller and about 30\% faster, but it does not survive
post-training INT8 quantisation: its field macro F1 fell from 0.834 (FP32) to
0.201 (INT8), a 63-point drop, while Large loses only 2.6 points. Large INT8
already runs in 70 to 300\,ms on the phone, so the speed gain does not justify the
accuracy loss. Quantisation-aware training could rescue Small; it is a Phase 2 option.

\subsection{Confidence threshold and out-of-scope photos}

A diagnosis is shown only when the model is confident and the photo looks like the
chosen crop. With temperature $T = 0.687$ and
$p = \mathrm{softmax}(\mathrm{logits}/T)$ over all 10 classes:

\begin{itemize}
  \item \textbf{crop mass} $=$ sum of $p$ over the chosen crop's classes. If it is
        below \textbf{0.80}, the photo does not look like a leaf of that crop
        (``not sure, check the crop you chose'').
  \item \textbf{confidence} $=$ top class of the chosen crop, divided by the crop
        mass. If it is below \textbf{0.56}, the app says ``not sure''.
  \item Either case, or a failed quality gate, means no disease is named, no
        treatment is given, and the farmer is referred to the KVK.
\end{itemize}

Both values were chosen on the validation set with a stated rule: at least 97\% of
answered photos correct and at least 94\% of photos answered, then as few
out-of-scope answers as possible. Out-of-scope photos were 325 leaves of other
crops (each scored as tomato and as potato).

\begin{table}[H]
\centering
\begin{tabular}{@{}lrrr@{}}
\toprule
\hdr{Set} & \hdr{Answered} & \hdr{Answered and correct} & \hdr{Other-crop leaves answered} \\
\midrule
Validation & 94.4\% & 97.5\% & 19.7\% \\
Field test & 93.6\% & 95.8\% & --- \\
\bottomrule
\end{tabular}
\end{table}

A plain confidence cut-off of 0.52 would answer more photos but would give a
confident answer to 61\% of the other-crop leaves; the crop-mass check is what
removes most of them. About one in five other-crop leaves still gets an answer
(29\% when tomato is chosen, 11\% for potato), so the ``not tomato or potato''
rejection is a safeguard, not a guarantee.

% --------------------------------------------------------------- severity ---
\section{Severity}

The leaf and the lesions are segmented by colour (HSV thresholds with
morphological clean-up), and severity is lesion area divided by leaf area. It is a
transparent baseline, not a trained segmenter: on PlantSeg's expert masks the
lesion IoU is low (mean 0.19), and virus classes, which show as yellowing rather
than spots, are the weakest. The bands should be read as a hint, and the advice
says so. The Kotlin port matches the Python reference on all 8 golden photos to
0.1 points and takes 12 to 42\,ms per photo on the phone after the first.

A trained lesion segmenter on PlantSeg's masks is the planned upgrade.

% --------------------------------------------------------------- advisory ---
\section{Advisory}

\subsection{Knowledge base}

For each class: cause, conditions that favour it, symptoms, cultural control, and
approved chemical options with dose per acre, water volume and pre-harvest
interval. Sources: NIPHM / Directorate of PPQS AESA-based IPM packages for tomato
and potato, and the TNAU Agritech portal. Chemicals banned or restricted in India
are on a deny list.

\subsection{Language model}

Llama 3.2 1B Instruct, SpinQuant INT4 (1.1\,GB), run with ExecuTorch on the phone.
The prompt carries fixed rules plus only the KB facts for this case and the
farmer's answers. Greedy decoding. English uses a 768-token context; Hindi uses
1,024 because Hindi needs about $3\times$ the tokens.

\subsection{Safety check (every answer)}

\begin{itemize}
  \item No chemical, dose or number that is not in the KB facts for this case.
  \item No banned chemical.
  \item KVK / Kisan Call Centre referral present for severe or uncertain cases.
  \item Answer not cut off, and in the right script for the language.
\end{itemize}

If the only failure is a missing referral, the app adds the fixed referral
sentence and checks again. Otherwise the KB template is shown. Uncertain results
never go to Llama.

\subsection{Evaluation (60 cases: 10 classes $\times$ severities $\times$ English and Hindi)}

\begin{table}[H]
\centering
\begin{tabularx}{\textwidth}{@{}L l@{}}
\toprule
\hdr{Measure} & \hdr{Result} \\
\midrule
Llama output grounded (no chemical, dose or number outside the KB) & 52/52 (100\%) \\
Llama output passed every check & 39/52 (75\%); English 88\%, Hindi 62\% \\
Llama text shown after the referral retry & 40/52 (77\%) \\
Advice shown to the farmer that passed every check & \textbf{60/60 (100\%)} \\
\bottomrule
\end{tabularx}
\end{table}

Most failures are answers cut off at the token limit, mainly in Hindi. The
fallback to the template means the farmer never sees a failed answer.

\subsection{Languages}

English and Hindi are written by Llama (Hindi KB text pending review by a native
speaker). Bengali, Marathi, Tamil and Telugu use translated KB templates, marked
as not yet reviewed.

% ----------------------------------------------------- on-device results ----
\section{On-device performance (Redmi 12 5G)}

Snapdragon 4 Gen 2 (Arm Cortex-A78 and A55), Android 15, about 1.8 to 2.3\,GB free
RAM. ExecuTorch 1.5.1.

\begin{table}[H]
\centering
\begin{tabularx}{\textwidth}{@{}L l@{}}
\toprule
\hdr{Stage} & \hdr{Time} \\
\midrule
Leaf classifier & 70 to 300\,ms (first run slowest) \\
Quality gate & included above \\
Severity & 12 to 42\,ms per photo after the first (about 2\,s on the first photo) \\
Llama load & 2.4 to 7.4\,s \\
Llama first token, full grounded prompt & about 8\,s English (230 prompt tokens), about 19\,s Hindi \\
Llama decode & 10 to 14 tokens/s \\
Full flow, photo to Llama advice (airplane mode, 10 runs) & 18 to 19\,s \\
\bottomrule
\end{tabularx}
\end{table}

\begin{table}[H]
\centering
\begin{tabularx}{\textwidth}{@{}L l@{}}
\toprule
\hdr{Memory and energy} & \hdr{Value} \\
\midrule
Peak app RSS with Llama loaded & 1.51 to 1.62\,GB \\
Llama only started if free RAM $\geq$ 1,900\,MB; otherwise template advice & --- \\
Energy per diagnosis (Android batterystats, CPU only, screen excluded) & about 3.2\,mAh \\
Diagnoses per full charge (4,381\,mAh) & about 1,360 \\
\bottomrule
\end{tabularx}
\end{table}

The template advice appears at once; the Llama text replaces it only if it passes
the safety check, so the farmer is never waiting on a blank screen.

\textbf{Offline proof.} Airplane mode on, no active network, ping fails; 10 of 10
full flows completed with Llama advice (\texttt{fails=[]}). Shown in the
\href{https://drive.google.com/file/d/1YLU4_h4lRAjxv5cIshW1-GhEhf6i896o/view?usp=drive_link}{demo video}.

% ------------------------------------------------------------ limitations ---
\section{Limitations}

\begin{enumerate}
  \item \textbf{Accuracy on photos from new farms is lower} than the headline:
        about 70 to 80\% on the web-scraped sources. Our own phone photos from
        Indian fields would be the most valuable addition.
  \item \textbf{Close-up leaf photos work best.} The model sees the centre of the
        photo and gives one label. Whole-plant photos, or two problems in one
        photo, are not handled. A leaf or lesion detector is the next step.
  \item \textbf{Severity is a colour baseline}, not a trained segmenter.
  \item \textbf{Out-of-scope rejection is partial.} About one in five leaves of
        other crops still gets a confident answer. In a small extra check, 4 of 15
        photos of everyday objects got a confident answer when tomato was chosen
        (none when potato was chosen), and 24 of 60 photos of tomato diseases
        outside our list (bacterial spot, target spot, mosaic virus, spider mites)
        were given one of our labels. The app's photo hint and the KVK referral
        reduce the risk, and adding these conditions as classes, or an ``other''
        class, is planned.
  \item \textbf{Translations are not yet reviewed} by native speakers.
  \item \textbf{Energy is a software estimate}, not a meter reading, and excludes
        the screen.
  \item Some field classes have few test images (21 to 30 for four tomato
        classes), so their scores are noisy.
\end{enumerate}

% --------------------------------------------------------------- how to run ---
\section{How to run}

See the README: build the APK with Gradle, install with \texttt{adb install}, and
push \texttt{llama.pte} and \texttt{tokenizer.model} to the app's files folder.
The leaf model is bundled in the APK.

% ------------------------------------------------------------- references ---
\section{References}

\subsection*{Datasets}

\begin{itemize}[label={},leftmargin=2.4em,labelsep=0.45em,itemsep=0.3em]
  \item[{[1]}] Hughes, D.\,P.\ and Salathe, M. ``An open access repository of images on plant health to enable the development of mobile disease diagnostics.'' arXiv:1511.08060, 2015. \url{https://github.com/spMohanty/PlantVillage-Dataset}
  \item[{[2]}] Singh, D.\ et al.\ ``PlantDoc: A Dataset for Visual Plant Disease Detection.'' CoDS-COMAD 2020. \url{https://github.com/pratikkayal/PlantDoc-Dataset}
  \item[{[3]}] Wei, T.\ et al.\ ``PlantSeg: A Large-Scale In-the-wild Dataset for Plant Disease Segmentation.'' arXiv:2409.04038, 2024. \url{https://zenodo.org/records/17719108}
  \item[{[4]}] Gehlot, M., Saxena, R.\,K.\ and Gandhi, G.\,C.\ ``Tomato-Village: a dataset for end-to-end tomato disease detection in a real-world environment.'' Multimedia Systems, 2023. \url{https://github.com/mamta-joshi-gehlot/Tomato-Village}
  \item[{[5]}] Shabrina, N.\,H.\ et al.\ ``A novel dataset of potato leaf disease in uncontrolled environment.'' Data in Brief 52, 109955, 2024. \url{https://data.mendeley.com/datasets/ptz377bwb8/1}
  \item[{[6]}] ``Potato Leaf (Healthy and Late Blight).'' Mendeley Data, doi:10.17632/v4w72bsts5.2. \url{https://data.mendeley.com/datasets/v4w72bsts5/2}
  \item[{[7]}] Rashid, J.\ et al.\ ``Multi-Level Deep Learning Model for Potato Leaf Disease Recognition.'' Electronics 10(17), 2064, 2021. \url{https://www.kaggle.com/datasets/rizwan123456789/potato-disease-leaf-datasetpld}
\end{itemize}

\subsection*{Models and software}

\begin{itemize}[label={},leftmargin=2.4em,labelsep=0.45em,itemsep=0.3em]
  \item[{[8]}] Howard, A.\ et al.\ ``Searching for MobileNetV3.'' ICCV 2019. \url{https://arxiv.org/abs/1905.02244}
  \item[{[9]}] Meta. Llama 3.2 1B Instruct, SpinQuant INT4 ExecuTorch build. \url{https://huggingface.co/executorch-community/Llama-3.2-1B-Instruct-SpinQuant_INT4_EO8-ET} (Llama 3.2 Community License; Acceptable Use Policy: \url{https://www.llama.com/llama3_2/use-policy})
  \item[{[10]}] Liu, Z.\ et al.\ ``SpinQuant: LLM Quantization with Learned Rotations.'' arXiv:2405.16406, 2024. \url{https://arxiv.org/abs/2405.16406}
  \item[{[11]}] PyTorch ExecuTorch, on-device inference runtime. \url{https://github.com/pytorch/executorch}
  \item[{[12]}] Guo, C.\ et al.\ ``On Calibration of Modern Neural Networks'' (temperature scaling). ICML 2017. \url{https://arxiv.org/abs/1706.04599}
\end{itemize}

\subsection*{Crop protection knowledge}

\begin{itemize}[label={},leftmargin=2.4em,labelsep=0.45em,itemsep=0.3em]
  \item[{[13]}] NIPHM and Directorate of PPQS, Ministry of Agriculture and Farmers Welfare. ``AESA based IPM Package: Tomato.'' \url{https://niphm.gov.in/IPMPackages/Tomato.pdf}
  \item[{[14]}] NIPHM and Directorate of PPQS, Ministry of Agriculture and Farmers Welfare. ``AESA based IPM Package: Potato.'' \url{https://niphm.gov.in/IPMPackages/Potato.pdf}
  \item[{[15]}] TNAU Agritech Portal, crop protection pages for tomato and potato. \url{https://agritech.tnau.ac.in}
  \item[{[16]}] Cornell Vegetables, ``Tomato Leaf Mold.'' \url{https://vegetables.cornell.edu/crops/tomatoes/tomato-leaf-mold}
  \item[{[17]}] Royal Horticultural Society, ``Tomato leaf mould.'' \url{https://www.rhs.org.uk/disease/tomato-leaf-mould}
\end{itemize}

\subsection*{Challenge}

\begin{itemize}[label={},leftmargin=2.4em,labelsep=0.45em,itemsep=0.3em]
  \item[{[18]}] Arm Education. ``Bharat AI-SoC Challenge 2026--27.'' \url{https://arm-education.github.io/Arm-Developer-Labs/Bharat_AI_SoC_2026_27.html}
\end{itemize}

\vspace{0.8em}
\noindent\small Licences for every source are listed in \texttt{NOTICE.md} in the repository.

\end{document}
